\newcommand{\resolved}{{\color{cern@blue}[RESOLVED]}}
\chapter{Questions and Issues}
\label{questions_and_issues}

\section{Success and Failure for Archive Messages}

\textit{Results of a discussion between Jozsef, Giuseppe and Eric about the success and failure behaviour for archive requests:}

The current behaviour in CASTOR is that the file remains to be migrated until all copies have been successfully
migrated. Failed migration jobs get deleted, so the file cannot be \texttt{stager\_rm}ed or garbage collected
before an operator intervenes.

We think a similar behaviour should be implemented in EOSCTA:

\begin{itemize}
   \item The file will be opened to garbage collection when all copies are successfully archived to tape.
      CTA will not report success before that (this is already the case).
   \item When a failure occurs (after exhausting retries), CTA will report the error to EOS, with a new
      error-reporting URL (to be implemented). The job will then be placed in a failed job queue to be
      handled by operators.
   \item EOS will keep track of and expose to the user only the latest error (we have potentially one
      per tape copy, and if operator decides to retry the job entirely, the error could be reported again).
   \item EOS will clear the errors when receiving a success.
\end{itemize}

\section{Immutable files}

The files with an archive on tape should be immutable in EOS (raw data use case), or a delayed archive mechanism
should be devised for mutable files (CERNBox archive use case).

Immutability of a file is guaranteed by adding \texttt{u!} to the EOS ACL.

Currently we do not enforce this on the CTA Frontend side, we just assume EOS is taking care of it.

If we decide it's useful for CTA to check immutability of archived files, we could send the ACL across with the xattrs.
This is not sent at the moment, because all system and user attributes are filtered out.

\section{When can files be deleted?}

Disk copies cannot be deleted before they are archived on tape (pinning).

The full file could still be deleted, potentially leading to issues to be handled in the tape archive session.

\section{What should be the protocol for fast reconciliation?}

The workflow will both trigger the synchronous archive queuing and post a second delayed workflow job that will
check and re-issue the request if needed (in case the request gets lost in CTA). This event-driven reconciliation acts
as a fast reconciliation. The criteria to check the file status will be the EOS side status which CTA reports
asynchronously to EOS (see~\S\ref{dataSerialization}).

\section{When a file has multiple tape copies, when are notifications are sent to EOS?}

EOS will need to represent and handle part the tape status of the files. This includes the fact that the file should be
on tape, the name of the CTA storage class, and the mutually exclusive statuses indicated by CTA: not on tape, partially
on tape, fully on tape.	The report from CTA will use the ``tape replica'' message (see~\S\ref{dataSerialization}).

For CASTOR, there is an additional constraint that the disk copy cannot be deleted until all tape copies have been
successfully written. The above scheme keeps track of the number of tape copies written and it will be up to the
EOS developers to ensure that this constraint is observed.

In CASTOR, the following notifications are sent during archiving a file with $n$ tape copies:
\begin{itemize}
   \item On successful write of the first tape copy, the \textbf{m-bit} is set. This indicates to the experiment that
      they can safely delete their copy of the data.
   \item On successful write of the $n^{th}$ tape copy, the \textbf{CAN\_BE\_MIGR} status is set in the database. This
      indicates that the file can be deleted from CASTOR's staging area.
\end{itemize}

For CTA, at what point(s) should we notify EOS that a file has been archived?
\begin{itemize}
   \item After the first copy is archived?
   \item After each copy is archived?
   \item After the $n^{th}$ copy is archived?
\end{itemize}

\gitlab{228}{Test archving files with \texttt{num\_copies} $gt 1$}

\section{Should the CTA catalogue methods prepareForNewFile() and prepareToRetrieveFile() detect repeated requests from
EOS instances?}

EOS does not keep track of requests which have been issued. We have said that CTA should implement idempotent retrieve queuing.

What are the consequences if we do not implement idempotent retrieve queuing?

What about archives and deletes?

\subsection{If so how should the catalogue communicate such ``duplicate'' requests to the caller (Scheduler\slash cta-frontend plugin)?}

The CTA Frontend calls the Scheduler which calls the Catalogue.
There are several possible schemes for handling duplicate jobs:
\begin{enumerate}
   \item If duplicates are rare, perhaps they don't need to be explicitly handled
   \item When a retrieve job is submitted, the Scheduler could check in the Catalogue for duplicates 
   \item When a retrieve job completes, the Tape Server could notify the Scheduler, which could then check for and
      drop any duplicate jobs in its queue.
\end{enumerate}

Reporting of retrieve status could set an \texttt{xattr}. Then the user would be able to monitor status which could
reduce duplicate requests.

Failed archivals or other CTA errors could also be logged as an \texttt{xattr}.

\subsection{If the CTA catalogue keeps an index of ongoing archive and retrieve requests, what will be the new
protocol additions (EOS, cta-frontend and cta-taped) required to guarantee that ``never completed'' requests are removed
from the catalogue?}

Such a protocol addition could be something as simple as a timeout.
                     
% \section{How do we deal with the fact that the current C++ code of the EOS/CTA interface that needs to be compiled on
% the EOS side on SLC6 will not compile because it uses std::future?}
%
% Please could you take on the responsibility of addressing the EOS/CTA interface issues described by these questions.
% You do NOT need to do any work towards these issues before the next Wednesday meeting.  These issues need to be
% vetted by the whole team during the meeting.  What's left is then under your responsibility.
% I am sending you these questions in order to start the process of concretely describing the scope of your EOS/CTA
% interface work. More issues are described within the EOS/CTA interface document of Eric.
% 
% If and only if there is time during the next tape developments meeting, we should try to vet the above questions and
% decide on the most appropriate place to have them written up, in  other words do they become gitlab issues or do they
% get added to the EOS/CTA interface document?  If we don’t have time during the meeting then those interested can have
% a separate meeting sometime later.


\section{CTA Failure}

What is the mechanism for restarting a failed archive request (in the case that EOS accepts the request and CTA fails
subsequently)?

If CTA is unavailable or unable to perform an archive operation, should EOS refuse the archive request and report failure
to the User?

What is the retry policy?

\section{File life cycle}

Full life cycle of files in EOS with copies on tape should be determined (they inherit their tape properties
from the directory, but what happens when the file gets moved or the directory properties changed?).

\section{Storage Classes}

The list of valid storage classes needs to be synchronized between EOS and CTA. EOS should not allow a power user to
label a directory with an invalid storage class. CTA should not delete or invalidate a storage class that is being used
by EOS.

\section{Request Queue}

Chaining of archive and retrieve requests to retrieve requests.

Execution of retrieve requests as disk to disk copy if possible.

%Catalogue will also keep track of requests for each files (archive and retrieve) so that queueing can be made idempotent.

\section{Catalogue}

Catalogue files could hold the necessary info to recreate the archive request if needed.

\section{Questions administrators need to be able to answer}

The \texttt{cta-admin} command should be include functions to allow administrators to answer the following questions:

\begin{itemize}
\item Why is data not going to tape?
\item Why is data not coming out of tapes?
\item Which user is responsible for system overload?
\end{itemize}

\section{Return value \resolved}

Notification return structure for synchronous workflows contains the following:
\begin{itemize}
   \item Success code (\texttt{RSP\_SUCCESS})
   \item A list of extended attributes to set (\textit{e.g.}, set the ``CTA archive ID'' \texttt{xattr} of the EOS file being queued for archival)
   \item Failure code (\texttt{RSP\_ERR\_PROTOBUF}, \texttt{RSP\_ERR\_CTA} or \texttt{RSP\_ERR\_USER})
   \item Failure message which can be logged by EOS or communicated to the end user (\textit{e.g.}, ``Cannot open file for writing because there is no route to tape'')
\end{itemize}

\begin{lstlisting}
message Response {
  enum ResponseType {
    RSP_INVALID             = 0;    //< Response type was not set
    RSP_SUCCESS             = 1;    //< Request is valid and was accepted for processing
    RSP_ERR_PROTOBUF        = 2;    //< Framework error caused by Google Protocol Buffers layer
    RSP_ERR_CTA             = 3;    //< Server error reported by CTA Frontend
    RSP_ERR_USER            = 4;    //< User request is invalid
  }
  ResponseType type         = 1;    //< Encode the type of this response
  map<string, string> xattr = 2;    //< xattribute map
  string message_txt        = 3;    //< Optional response message text
}
\end{lstlisting}

\section{Will EOS instance names within the CTA catalogue be ``long'' or ``short''? \resolved}

\textit{We all agreed to use ``long'' EOS instance names within CTA and specifically the CTA catalogue. An example of a
long EOS instance name is ``eosdev'' with its corresponding short instance name being ``dev''.}
\begin{flushright}
--- Minutes from today's tape developments meeting, Wed 22 Nov 2017
\end{flushright}

This implies that there will be a separate instance name for each VO (``eosatlas'', ``eoscms'', \textit{etc.}) and a
unique SSS key for each instance name.

\section{Do we want the EOS namespace to store CTA archive IDs or not? \resolved}

\begin{description}
   \item[If no:] we are allowing that the EOS file ID uniquely identifies the file. We must maintain a one-to-one mapping
      from EOS ID to CTA archive ID on our side. This also implies that the file is immutable.
   \item[If yes:] we must generate the CTA archive ID and return it to EOS. There must be a guarantee that EOS has attached
      the archive ID to the file (probably as an xattr but that's up to the EOS team), i.e. \textbf{the EOS end-user must
      never see an EOS file with a tape replica but without an archive ID}. EOS must provide the CTA archive ID as the
      key to all requests. 
\end{description}

\subsection*{Solution}

Archive IDs will be allocated by CTA when a file is created. The Archive ID will be stored in the EOS
namespace as an extended attribute of the file. EOS must use the archive ID to archive, retrieve or delete files.

Archive IDs are not file IDs, \textit{i.e.} the archive ID identifies the version of the file that was archived. In the
case of Physics data, the files should be immutable so in practice there is one Archive ID per file.

In the backup use case, if we allowed mutable files, we would need a mechanism to track archived file versions. On the
EOS side, changes to files are versioned, so each time a file is updated, the Archive ID should also be updated.
Old versions of the file would maintain a link to their archive copy via the versioned extended attributes. But in this
case we probably also need a way to mark archive copies of redundant versions of files for deletion.

\subsection*{Design notes from Steve}

\textit{One of the reasons I wanted an archive ID in the EOS namespace was that I wanted to have one primary key for the CTA
file catalogue and I wanted it to be the CTA archive ID. Therefore I expected that retrieve and delete requests issued
by EOS would use that key.}

\textit{This ``primary key'' requirement is blown apart by the requirement of the CTA catalogue to
identify duplicate archive requests.  The CTA archive ID represents an ``archive request'' and not an individual EOS file.
Today, 5 requests from EOS to archive the same EOS file will result in 5 unique CTA archive IDs. Making the CTA catalogue
detect 4 of these requests as duplicate means adding a ``second'' primary key composed of the EOS instance name and the EOS
file ID. It also adds the necessity to make sure that archive requests complete in the event of failure, so that retries
from EOS will eventually be accepted and not forever refused as duplicate requests.  It goes without saying that dropping
the CTA archive ID from EOS also means using the EOS instance name and EOS file ID as primary key for retrieve and delete
requests from EOS.}

\textit{The requirement for a ``second'' primary key may be inevitable for reasons other than (idempotent) archive, retrieve and
delete requests from EOS.  CTA tape operators will want to drill down into the CTA catalogue for individual end user files
when data has been lost or something has ``gone wrong''.  The question here is, should it be a ``primary key'' as in no
duplicate values or should it just be an index for efficient lookup?}

